\documentclass[11pt]{letter}
\usepackage[margin=1in]{geometry}
\signature{Krishna Harish\\Elkins High School, Missouri City, TX\\krishnaharish2009@gmail.com}
\address{Krishna Harish\\Elkins High School\\Missouri City, TX, USA}
\begin{document}
\begin{letter}{The Editors\\Experimental Mathematics}

\opening{Dear Editors,}

Please consider the enclosed manuscript, ``Directed Weisfeiler--Leman
refinement versus the invariants of the magnitude--path spectral
sequence,'' for publication in Experimental Mathematics.

The paper is a computer-assisted study of how much two families of
digraph invariants can tell apart: the homology theories organized by
the magnitude--path spectral sequence (magnitude, path, and
reachability homology), and directed Weisfeiler--Leman color
refinement, the standard bound on the separating power of
message-passing networks on digraphs. Working from an exhaustive exact
enumeration of all 1{,}550{,}786 digraphs on up to six vertices, the
computation locates every minimal witness that distinguishes the two
families, and this experimental picture then drives the theory: the
families are incomparable, with minimal witnesses pinned down exactly;
the bigraded rank function of the first page of the spectral sequence
determines neither its second page nor its target; and a
component-counting description of degree-two magnitude homology,
verified against the full enumeration, supplies closed formulas that
explain each separation.

The work is experimental in method and exact in execution. All homology
ranks are computed over the rationals with no floating point; two
independent implementations sharing no linear algebra code
cross-validate each other; the degree-two formulas were checked against
the general implementation on all 1{,}550{,}786 digraphs; and every
reported witness was re-verified from scratch. The complete toolkit and
data, with scripts that regenerate every number and table, are public
on GitHub and archived at Zenodo (DOI~10.5281/zenodo.21138733), so every
computational claim in the paper is independently reproducible.

To my knowledge this is the first software to compute magnitude, path,
and reachability homology of general digraphs in a single
exact-arithmetic framework, and the first placement of these invariants
relative to directed refinement; the undirected analogue, persistent
homology against the Weisfeiler--Leman hierarchy, was carried out by
Ballester and Rieck (2024).

This work is not under consideration elsewhere, and I have no competing
interests.

\closing{Sincerely,}
\end{letter}
\end{document}
